\documentclass[11pt,letterpaper]{article}
\usepackage[utf8]{inputenc}
\usepackage[provide=*,english]{babel}
\usepackage{amsmath,amssymb,amsfonts}
\usepackage[margin=2.5cm]{geometry}
\usepackage{fancyhdr}
\usepackage{enumitem}
\usepackage[most]{tcolorbox}
\usepackage{xcolor}
\usepackage{tikz}
\usetikzlibrary{arrows.meta,angles,quotes,calc,babel}
\usepackage{graphicx}
\usepackage{needspace}

% Palette and layout of the reference document.
\definecolor{blackmain}{HTML}{1A1A1A}
\definecolor{orangeacc}{HTML}{E65C00}
\definecolor{creambg}{HTML}{FBF8F1}
\definecolor{orangebg}{HTML}{FFF4EB}
\definecolor{graymuted}{HTML}{7A746E}
\definecolor{pureblack}{HTML}{0D0D0D}

\setlength{\headheight}{24pt}
\pagestyle{fancy}
\fancyhf{}
\fancyhead[L]{\textcolor{blackmain}{\textbf{\small Mathematical Methods I}}}
\fancyhead[R]{\textcolor{orangeacc}{\textbf{\small The metric and tensors}}}
\fancyfoot[C]{\textcolor{graymuted}{\thepage}}
\renewcommand{\headrulewidth}{0.8pt}
\renewcommand{\headrule}{\hbox to\headwidth{\color{orangeacc}\leaders\hrule height \headrulewidth\hfill}}

\setlist[itemize]{label={\small\textcolor{orangeacc}{$\blacksquare$}},leftmargin=1.6em}
\setlist[enumerate]{label={\textcolor{orangeacc}{\arabic*.}},leftmargin=1.8em}

\newcommand{\manualsection}[1]{%
  \Needspace{13\baselineskip}%
  \vspace{0.35cm}
  \par\noindent{\Large\bfseries\color{blackmain} #1}\par\nobreak
  \noindent{\color{orangeacc}\rule{\linewidth}{1.4pt}}\par\nobreak
  \vspace{0.15cm}
}
\newcommand{\manualsubsection}[1]{%
  \par\vspace{0.2cm}\noindent{\bfseries\color{orangeacc} #1}\par\nobreak
  \vspace{0.12cm}
}
\newtcolorbox{infobox}[2][orangeacc]{
  enhanced,colback=creambg,colframe=#1!90!black,
  boxrule=1pt,arc=2mm,left=9pt,right=9pt,top=5pt,bottom=5pt,
  before skip=7pt,after skip=7pt,
  before upper={\textbf{\color{#1!95!black}\large #2}\par\nobreak
    {\color{#1!50}\rule{\linewidth}{0.6pt}}\par\nobreak\smallskip}
}
\newtcolorbox{examplebox}[1]{
  enhanced,colback=orangebg,colframe=orangeacc!85!black,
  boxrule=1pt,arc=2mm,left=9pt,right=9pt,top=5pt,bottom=5pt,
  before skip=7pt,after skip=7pt,
  before upper={\textbf{\color{orangeacc!95!black}\large #1}\par\nobreak
    {\color{orangeacc!40}\rule{\linewidth}{0.6pt}}\par\nobreak\smallskip}
}
\newtcolorbox{theorembox}[1]{
  enhanced,colback=creambg!70!white,colframe=pureblack,
  boxrule=1.1pt,arc=2mm,left=9pt,right=9pt,top=5pt,bottom=5pt,
  before skip=7pt,after skip=7pt,
  before upper={\textbf{\color{pureblack}\large #1}\par\nobreak
    {\color{orangeacc}\rule{\linewidth}{0.8pt}}\par\nobreak\smallskip}
}
\newcommand{\term}[1]{\textbf{\color{orangeacc}#1}}
\renewcommand{\familydefault}{\rmdefault}
\tikzset{every node/.append style={font=\rmfamily}}

\begin{document}
\setlength{\abovedisplayskip}{6pt}
\setlength{\belowdisplayskip}{6pt}
\setlength{\abovedisplayshortskip}{4pt}
\setlength{\belowdisplayshortskip}{4pt}
\setlength{\jot}{2pt}
\color{blackmain}
\begin{center}
{\LARGE\textbf{Mathematical Methods I}}\\[0.2cm]
{\Large\textcolor{orangeacc}{\textbf{The metric and tensors}}}\\[0.4cm]
{\large\textcolor{graymuted}{Gabriel Velázquez}}\\[0.25cm]
\end{center}

\manualsection{1. Tensors and their components}
A tensor generalizes the notions of a scalar and a vector. In a coordinate system, its components are identified by indices; the total number of indices is its \term{rank}.
\begin{infobox}[orangeacc]{Types of tensor}
\begin{itemize}[topsep=2pt,itemsep=3pt]
\item \textbf{Rank 0:} a scalar, such as mass or temperature at a point.
\item \textbf{Rank 1:} a vector, such as velocity or force, or a covector, such as the differential of a function.
\item \textbf{Rank 2 or higher:} a tensor whose components form a matrix or a higher dimensional array.
\end{itemize}
In dimension $d$, a rank $n$ tensor has $d^n$ components in a basis, although symmetries may reduce the number of independent components.
\end{infobox}
The defining feature of a tensor is its \term{transformation law}: its components change in a prescribed way when coordinates change, while the geometric object they represent remains the same.
\begin{theorembox}{Geometric definition}
At a point $p$, a tensor of type $(r,s)$ is a multilinear map
\[
T_p:\underbrace{T_p^*M\times\cdots\times T_p^*M}_{r\text{ covectors}}
\times\underbrace{T_pM\times\cdots\times T_pM}_{s\text{ vectors}}
\longrightarrow\mathbb R.
\]
Its components carry $r$ upper indices and $s$ lower indices. For example, $A^i$ is contravariant, $A_i$ is covariant, and $T^i{}_j$ is mixed.
\end{theorembox}

\newpage
\manualsection{2. The metric: from vectors to covectors}
Tangent vectors and covectors belong to different spaces. A \term{metric} $g$ assigns a real number to each pair of tangent vectors, bilinearly and symmetrically. Because it is nondegenerate, it associates each vector $v$ with the covector $v^\flat=g(v,\,\cdot\,)$; this operation is called \term{lowering an index}.
\begin{infobox}[orangeacc]{Covariant metric tensor}
In coordinates $x^1,\ldots,x^d$, let $e_i=\partial/\partial x^i$ be the tangent basis. The metric components are
\[
\boxed{g_{ij}=g(e_i,e_j)=\left\langle\frac{\partial}{\partial x^i},
\frac{\partial}{\partial x^j}\right\rangle,\qquad g_{ij}=g_{ji}.}
\]
The matrix $(g_{ij})$ represents a symmetric tensor of type $(0,2)$. In the Euclidean or Riemannian case, $g(v,v)>0$ whenever $v\ne0$.
\end{infobox}
The components $v_i=g_{ij}v^j$ are those of $v^\flat$; here and below, repeated indices are summed. The \term{inverse metric} has components $g^{ij}$ and satisfies
\[
\boxed{g^{ik}g_{kj}=\delta^i{}_j,\qquad v^i=g^{ij}v_j.}
\]
Thus $g^{ij}$ is symmetric and of type $(2,0)$, and it raises an index. The Kronecker delta $\delta^i{}_j$ is the identity matrix.

\manualsection{3. What $ds^2$ represents}
The symbol $ds^2$ denotes the \term{line element}: the metric evaluated on an infinitesimal displacement. In Riemannian geometry, it is the square of the infinitesimal distance between nearby points. It is also called the \term{first fundamental form}.
\begin{theorembox}{First fundamental form}
\[
\boxed{ds^2=g_{ij}\,dx^i\,dx^j.}
\]
The Einstein convention is used: repeated indices $i$ and $j$ are summed. The products $dx^i\,dx^j$ express the quadratic evaluation of the metric on a displacement; they are not exterior products $dx^i\wedge dx^j$.
\end{theorembox}
\manualsubsection{The metric in action}
For tangent vectors $v=v^ie_i$ and $w=w^ie_i$, the inner product and Riemannian norm are
\[
\langle v,w\rangle=g_{ij}v^iw^j,
\qquad \lVert v\rVert=\sqrt{g_{ij}v^iv^j}.
\]
For nonzero vectors, the angle $\phi$ is determined by
\[
\cos\phi=\frac{g_{ij}v^iw^j}{\sqrt{g_{ij}v^iv^j}\,\sqrt{g_{ij}w^iw^j}}.
\]
If $x^i(t)$ parametrizes a curve with a Riemannian metric, its length from $t=a$ to $t=b$ is obtained by integrating the speed measured by $g$:
\[
\boxed{L=\int_a^b\sqrt{g_{ij}(x(t))\frac{dx^i}{dt}\frac{dx^j}{dt}}\,dt.}
\]
In dimension $d$, the metric also determines the volume element $dV_g=\sqrt{\det(g_{ij})}\,dx^1\cdots dx^d$ in oriented Riemannian coordinates. It therefore enters the measurement of both curves and regions.

\manualsection{4. From the tensor product to the wedge product}
\begin{infobox}[orangeacc]{Tensor product and wedge product}
For covectors $\alpha$, $\beta$, the tensor product is defined by
\[
\boxed{(\alpha\otimes\beta)(v,w)=\alpha(v)\beta(w).}
\]
It need not be symmetric or alternating. In contrast,
\[
\boxed{\alpha\wedge\beta=\alpha\otimes\beta-\beta\otimes\alpha.}
\]
Thus $dx^1\otimes dx^1$ is nonzero, although $dx^1\wedge dx^1=0$.
\end{infobox}
The identity shows what changes when the two vectors are exchanged. Evaluating the products on $v,w$ gives
\begin{theorembox}{Evaluation and alternation}
\[
(\alpha\wedge\beta)(v,w)
=\alpha(v)\beta(w)-\alpha(w)\beta(v).
\]
Thus $(\alpha\wedge\beta)(w,v)=-(\alpha\wedge\beta)(v,w)$ and $\alpha\wedge\alpha=0$. In contrast, $(dx^1\otimes dx^1)(v,v)=[dx^1(v)]^2$ can be nonzero.
\end{theorembox}
\begin{examplebox}{Example in the plane}
If $\alpha=a\,dx+b\,dy$ and $\beta=c\,dx+d\,dy$, then
\[
\alpha\wedge\beta=(ad-bc)\,dx\wedge dy.
\]
In particular, for $u=(u^x,u^y)$ and $v=(v^x,v^y)$,
\[
(dx\wedge dy)(u,v)=u^xv^y-u^yv^x
=\det\begin{pmatrix}u^x&v^x\\u^y&v^y\end{pmatrix}.
\]
The wedge product records oriented area; $dx\otimes dy$ alone gives $u^xv^y$ and does not necessarily change sign when the vectors are exchanged.
\end{examplebox}
\manualsection{5. Example: the metric in polar coordinates}
In the Euclidean plane, let $u^1=r$, $u^2=\theta$, and
\[
x=r\cos\theta,\qquad y=r\sin\theta,\qquad r>0.
\]
To construct the metric, differentiate each Cartesian coordinate with respect to $r$ and $\theta$:
\begin{align*}
\frac{\partial x}{\partial r}&=\cos\theta,
&\frac{\partial y}{\partial r}&=\sin\theta,\\
\frac{\partial x}{\partial\theta}&=-r\sin\theta,
&\frac{\partial y}{\partial\theta}&=r\cos\theta.
\end{align*}
The associated tangent vectors are therefore
\[
e_r=(\cos\theta,\sin\theta),\qquad
e_\theta=(-r\sin\theta,r\cos\theta).
\]
The Cartesian metric on the plane is Euclidean. Hence,
\[
g_{ab}=e_a\cdot e_b
=\sum_{k=1}^{2}\frac{\partial x^k}{\partial u^a}
                  \frac{\partial x^k}{\partial u^b},
\qquad (x^1,x^2)=(x,y),\ (u^1,u^2)=(r,\theta).
\]
Compute its three independent components separately:
\begin{align*}
g_{rr}&=\left(\frac{\partial x}{\partial r}\right)^2
       +\left(\frac{\partial y}{\partial r}\right)^2
       =\cos^2\theta+\sin^2\theta=1,\\
g_{\theta\theta}&=\left(\frac{\partial x}{\partial\theta}\right)^2
       +\left(\frac{\partial y}{\partial\theta}\right)^2
       =r^2\sin^2\theta+r^2\cos^2\theta=r^2,\\
g_{r\theta}=g_{\theta r}&=
       \frac{\partial x}{\partial r}\frac{\partial x}{\partial\theta}
       +\frac{\partial y}{\partial r}\frac{\partial y}{\partial\theta}
       =-r\sin\theta\cos\theta+r\sin\theta\cos\theta=0.
\end{align*}
\begin{theorembox}{Metric matrix and line element in polar coordinates}
\[
\boxed{(g_{ab})=\begin{pmatrix}1&0\\0&r^2\end{pmatrix},
\qquad ds^2=dr^2+r^2d\theta^2.}
\]
The angular coordinate is converted to length by the scale factor $r$: at fixed radius, a change $d\theta$ corresponds to a distance $r\,d\theta$.
\end{theorembox}
\manualsection{6. Detailed calculation of a circle's circumference}
We seek the circumference of a circle of fixed radius $R>0$ using the polar metric $ds^2=dr^2+r^2d\theta^2$.
\begin{examplebox}{Step 1. Parametrize the curve}
Set $t=\theta$ and make one complete revolution for $0\leq t\leq2\pi$:
\[
u^1(t)=r(t)=R,\qquad u^2(t)=\theta(t)=t.
\]
The tangent vector therefore has components
\[
\frac{dr}{dt}=0,\qquad\frac{d\theta}{dt}=1.
\]
The first equation says that the curve has no radial motion; the second says that the angle increases by one unit per unit of $t$.
\end{examplebox}
\begin{examplebox}{Step 2. Evaluate the metric on the tangent}
In polar coordinates, $g_{rr}=1$, $g_{r\theta}=g_{\theta r}=0$, and $g_{\theta\theta}=r^2$. Expanding the Einstein sum gives
\begin{align*}
g_{ab}\frac{du^a}{dt}\frac{du^b}{dt}
&=g_{rr}\left(\frac{dr}{dt}\right)^2
  +2g_{r\theta}\frac{dr}{dt}\frac{d\theta}{dt}
  +g_{\theta\theta}\left(\frac{d\theta}{dt}\right)^2\\
&=(1)(0)^2+2(0)(0)(1)+(R^2)(1)^2=R^2.
\end{align*}
The speed measured by the metric is $\sqrt{R^2}=R$ because the radius is positive.
\end{examplebox}
\begin{examplebox}{Step 3. Integrate the length}
Substitute into the length formula, with limits $0$ and $2\pi$:
\[
\boxed{L=\int_0^{2\pi}\sqrt{g_{ab}\dot u^a\dot u^b}\,dt
=\int_0^{2\pi}R\,dt
=R[t]_{0}^{2\pi}=2\pi R.}
\]
This recovers the usual circumference from the line element in curvilinear coordinates.
\end{examplebox}

\newpage
\manualsection{7. Transformation and invariance}
The line element is the same geometric object in any coordinate system. Under a change from $x^i$ to coordinates $x'^a$, the differentials satisfy
\[
dx^i=\frac{\partial x^i}{\partial x'^a}\,dx'^a.
\]
Substituting them into $ds^2=g_{ij}dx^i dx^j$ gives the covariant transformation law:
\begin{theorembox}{Transformation law of the metric}
\[
\boxed{g'_{ab}=\frac{\partial x^i}{\partial x'^a}
\frac{\partial x^j}{\partial x'^b}\,g_{ij}.}
\]
Although the component matrix changes, $g'_{ab}dx'^a dx'^b=g_{ij}dx^i dx^j$.
\end{theorembox}
\begin{examplebox}{From Cartesian to polar coordinates}
In the plane, the Cartesian metric matrix is the identity. If $J=\partial(x,y)/\partial(r,\theta)$, then
\[
J=\begin{pmatrix}\cos\theta&-r\sin\theta\\
\sin\theta&r\cos\theta\end{pmatrix},\qquad
g_{\mathrm{polar}}=J^{\mathsf T}J
=\begin{pmatrix}1&0\\0&r^2\end{pmatrix}.
\]
This matrix multiplication reproduces the calculation of the products $e_a\cdot e_b$.
\end{examplebox}

\manualsection{8. Pythagoras in curvilinear coordinates}
In orthogonal coordinates $q^i$, the metric matrix is diagonal. For a Riemannian metric, define the \term{scale factor} $h_i=\sqrt{g_{ii}}$, with no sum over $i$. In three dimensions,
\begin{infobox}[orangeacc]{Orthogonal line element}
\[
\boxed{ds^2=(h_1\,dq^1)^2+(h_2\,dq^2)^2+(h_3\,dq^3)^2.}
\]
This is the Pythagorean theorem expressed through infinitesimal coordinate changes. In the polar plane, $h_r=1$ and $h_\theta=r$.
\end{infobox}

\newpage
\manualsection{9. The semi-Riemannian metric}
In Riemannian geometry, the metric is \term{positive definite}: $g(v,v)>0$ for every nonzero vector. This lets us interpret $ds^2$ as the square of an infinitesimal length. In \term{semi-Riemannian} geometry, the metric remains symmetric and nondegenerate but can have both positive and negative signs. Consequently, $g(v,v)$ may be positive, negative, or zero even when $v\ne0$.

The central physical example is a \term{Lorentzian} metric, with one time direction and three spatial directions. In flat Minkowski spacetime, using inertial coordinates and the sign convention $(-,+,+,+)$, the metric matrix is
\[
(g_{\mu\nu})=\operatorname{diag}(-c^2,1,1,1)
\quad\text{for coordinates }(t,x,y,z).
\]
\begin{theorembox}{Interval in Minkowski spacetime}
With the sign convention $(-,+,+,+)$ and inertial coordinates $(t,x,y,z)$,
\[
\boxed{ds^2=-c^2dt^2+dx^2+dy^2+dz^2.}
\]
Here $c$ is the speed of light. The value of the interval is preserved under coordinate changes.
\end{theorembox}

\end{document}
